% Source: https://artofproblemsolving.com/community/c3h1630145_math_problem_hard_4_me
% Stable identity: aops-topic-1630145-q01
The line $y=-2x+6$ meets the axes at $(3,0)$ and $(0,6)$; the line $x+y=2$ meets them at $(2,0)$ and $(0,2)$. The two lines meet where $-2x+6=2-x$, i.e. at $(4,-2)$, which is outside the first quadrant, so in the first quadrant the line $x+y=2$ lies entirely below $y=-2x+6$.

\begin{asy}
usepackage("amsmath");
size(6.5cm);
defaultpen(fontsize(10pt));
fill((2,0)--(3,0)--(0,6)--(0,2)--cycle,lightgray);
draw((-0.8,0)--(4.8,0),Arrow(TeXHead));draw((0,-2.6)--(0,7),Arrow(TeXHead));
label("$x$",(4.8,0),E);label("$y$",(0,7),N);
draw((-0.4,6.8)--(4.3,-2.6),blue+linewidth(0.9));draw((-0.6,2.6)--(4.8,-2.8),red+linewidth(0.9));
dot((2,0));dot((3,0));dot((0,2));dot((0,6));dot((4,-2));
label("$(2,0)$",(2,0),2*dir(235));label("$(3,0)$",(3,0),NE);label("$(0,2)$",(0,2),W);label("$(0,6)$",(0,6),W);label("$(4,-2)$",(4,-2),SW);
label("$y=-2x+6$",(1.3,3.6),E,blue);label("$x+y=2$",(4.8,-2.8),E,red);label("$7$",(0.9,2.2));
label("\mbox{}",(min(currentpicture,user=true)+max(currentpicture,user=true))/2,black); // keeps the text after the figure black when LaTeX typesets the labels inline
\end{asy}

The region between them is the large triangle with vertices $(0,0)$, $(3,0)$, $(0,6)$ with the small triangle $(0,0)$, $(2,0)$, $(0,2)$ removed:
\[
\dfrac{1}{2}\cdot3\cdot6-\dfrac{1}{2}\cdot2\cdot2=9-2=\boxed{7}.
\]
